\documentclass[11pt]{article}
\usepackage[margin=1in]{geometry}
\usepackage[T1]{fontenc}
\usepackage{amsmath,amssymb,amsthm,mathtools}
\usepackage[hidelinks]{hyperref}
\newtheorem{theorem}{Theorem}
\newtheorem{proposition}[theorem]{Proposition}
\newtheorem{lemma}[theorem]{Lemma}
\newtheorem{corollary}[theorem]{Corollary}
\newcommand{\Reg}{\operatorname{Reg}}
\newcommand{\Var}{\operatorname{Var}}
\newcommand{\KL}{\operatorname{KL}}
\newcommand{\ip}[2]{\langle #1,#2\rangle}
\newcommand{\norm}[1]{\lVert #1\rVert}
\title{Explicit special cases and a quantitative recurrence criterion for plain Optimistic Hedge}
\author{OpenAI GPT 6 Astra Pro}
\date{18 October 2026}
\begin{document}
\maketitle

\section*{Problem}
\textbf{prb_bc572aa70a8748a3b489fa9941f9307e: Explicit constant-regret bounds for plain Optimistic Hedge.}
The downloaded submission prompt identifies this entry as \texttt{ECON-2};
the permanent identifier above is used throughout this contribution.
\href{https://we-solve.github.io/problems/explicit-constant-regret-bounds-for-plain-optimistic-hedge/}{Public problem page}.

\section*{Title}
Explicit special cases and a quantitative recurrence criterion for plain Optimistic Hedge.

\section*{Claim type}
\textbf{Partial result, with a conditional general-game criterion. The full problem remains unresolved.}

\section*{Scope}
The exact algorithm in the problem admits: (i) the bound
$R_i(T)\leq4\log K$ at $\eta=\log2$ when player $i$ has a weakly dominant
pure action, even against arbitrary opponent play; (ii) the bound
$R_i(T)\leq(4N-2)\log K+3/2$ at
$\eta=1/[2(N-1)]$ in the class specified by
\eqref{eq:aggregate-condition}, which includes all two-player constant-sum
games; and (iii) an entirely numerical conversion from supplied finite
recurrence parameters to a horizon-uniform regret bound. A static-payoff
example also proves a lower bound $(\log K)/\eta-3/2$ for any proposed
step size. Section 6 conditionally computes general-game recurrence
parameters under decidability of the real exponential field. No
unconditional general-game parameter rule is supplied.

\section*{Assumptions and limitations}
There are $N\geq2$ players, $1\leq m_i\leq K$, and payoffs in $[0,1]$.
The feedback is the exact expected payoff vector and the regret is mixed-payoff
external regret against one fixed action. All logarithms are natural. No
sampled feedback, realized-action regret, floating-point implementation,
or last-iterate convergence is asserted. The special cases are included as
self-contained quantitative benchmarks, not claimed as new qualitative
resolutions of their previously studied classes. The recurrence criterion
is a quantitative development of Ha's method \cite{ha}. Section 6 gives
a terminating procedure to compute general-game bounds conditional on
decidability of the real exponential field. That additional assumption is
not proved here, and no running-time or polynomial bound is asserted.

\section*{Write-up}

\subsection*{1. Exact update and the outstanding question}
The update is precisely
\begin{equation}\label{eq:update}
 x_{ia}^{t+1}=\frac{x_{ia}^t\exp\{\eta(2v_{ia}^t-v_{ia}^{t-1})\}}
 {\sum_bx_{ib}^t\exp\{\eta(2v_{ib}^t-v_{ib}^{t-1})\}},
 \qquad x_{ia}^1=1/m_i,\quad v_i^0=0.
\end{equation}
Its telescoped representation is
\begin{equation}\label{eq:telescoped}
 x_{ia}^t\ \propto\
 \exp\!\left\{\eta\left(\sum_{s<t}v_{ia}^s+v_{ia}^{t-1}\right)\right\}.
\end{equation}
The exponent at $t=1$ is zero. At $t\geq2$, each earlier vector through
$t-2$ has coefficient one and the vector at $t-1$ has coefficient two.
This checks the initialization and the one-step optimistic term.

Ha's Theorem 4.1 asserts that for each action-count vector, some positive
step-size interval and some finite constant give
$R_i(T)\leq(\log m_i)/\eta+C_\star\eta$ uniformly in the game and horizon.
The text following Corollary 4.2 leaves dimension-explicit bounds open
\cite{ha}. Its finite-recurrence argument was inspected in Sections
5.3--5.5 and Appendices B--D. No contradiction to that argument was found
in this check. This is not a claim to have reconstructed every historical
analytic result on which its Noetherianity step relies.

\subsection*{2. Weakly dominant pure actions: a uniform explicit bound}

\begin{theorem}\label{thm:dominance}
Suppose player $i$ has a pure action $a_\star$ such that
$u_i(a_\star,a_{-i})\geq u_i(a,a_{-i})$ for every $a,a_{-i}$.
For every $\eta>0$, every opponent strategy sequence, and every $T$,
\[
 R_i(T)\leq\frac{e^\eta\log m_i}{1-e^{-\eta}}.
\]
Consequently, if each player has such an action, the common explicit choice
$\eta=\log2$ gives $R_i(T)\leq4\log K$ for every player and horizon.
No strictly positive dominance gap is required.
\end{theorem}

\begin{proof}
Suppress $i$. Put $d_a^t=v^t(a_\star)-v^t(a)$ and $d_a^0=0$.
Weak dominance and normalization imply $0\leq d_a^t\leq1$ and
$d_{a_\star}^t=0$. The action $a_\star$ maximizes cumulative payoff at
every horizon, so $R(T)=\sum_{t\leq T}\ip{x^t}{d^t}$.
Let $D_a^{t-1}=\sum_{s<t}d_a^s$, and define
\[
 W_t=\sum_a e^{-\eta D_a^{t-1}},\qquad
 y_a^t=\frac{e^{-\eta D_a^{t-1}}}{W_t}.
\]
Equation \eqref{eq:telescoped} gives
$x_a^t=y_a^t e^{-\eta d_a^{t-1}}/\sum_b y_b^t e^{-\eta d_b^{t-1}}$.
The denominator is at least $e^{-\eta}$ and the numerator factor at most
one, hence $x_a^t\leq e^\eta y_a^t$. Therefore
$\ip{x^t}{d^t}\leq e^\eta\ip{y^t}{d^t}$.

Convexity of $z\mapsto e^{-\eta z}$ on $[0,1]$ places its graph below its
endpoint chord: $e^{-\eta z}\leq1-(1-e^{-\eta})z$. Thus
\[
 \frac{W_{t+1}}{W_t}=\sum_a y_a^t e^{-\eta d_a^t}
 \leq1-(1-e^{-\eta})\ip{y^t}{d^t}.
\]
Since $-\log(1-z)\geq z$ for $0\leq z<1$,
\[
 (1-e^{-\eta})\sum_{t=1}^T\ip{y^t}{d^t}
 \leq\log W_1-\log W_{T+1}\leq\log m.
\]
The last inequality uses $W_1=m$ and the unit contribution of
$a_\star$ to $W_{T+1}$. Combining the two estimates proves the result,
including $m=1$, where the regret is zero.
\end{proof}

\subsection*{3. A general aggregate-comparator condition}
Suppose there is a fixed mixed-strategy profile $p=(p_1,\ldots,p_N)$
such that
\begin{equation}\label{eq:aggregate-condition}
 \sum_{i=1}^N\bigl[u_i(p_i,x_{-i})-u_i(x)\bigr]\geq0
 \qquad\text{for every mixed profile }x.
\end{equation}
The learners need not know $p$. Its existence is used only in the proof.
Condition \eqref{eq:aggregate-condition} is an additional game-class
assumption; it is not a property of every general-sum game.

\begin{theorem}\label{thm:aggregate}
Under \eqref{eq:aggregate-condition}, write
$L=\sum_j\log m_j$. If $0<\eta\leq1/[2(N-1)]$, the exact update
\eqref{eq:update} satisfies
\begin{equation}\label{eq:aggregate-bound}
 R_i(T)\leq\frac{\log m_i}{\eta}
 +\eta\bigl[1+4(N-1)L+4\eta^2N(N-1)\bigr]
 \quad\text{for every }i,T.
\end{equation}
In particular, $\eta=1/[2(N-1)]$ yields
$R_i(T)\leq(4N-2)\log K+3/2$.
\end{theorem}

\begin{proof}
For each player, define
$y_i^t\propto\exp(\eta\sum_{s<t}v_i^s)$ and set
\[
 a_{it}=\norm{y_i^{t+1}-x_i^t}_1,\qquad
 b_{it}=\norm{x_i^t-y_i^t}_1,\qquad
 d_{it}=\norm{v_i^t-v_i^{t-1}}_\infty.
\]
All coordinates of $x_i^t,y_i^t$ are positive at finite $t$.
For distributions $z,w$, let
\[
D(z\Vert w)=\sum_a z_a\log(z_a/w_a),
\]
with the usual $0\log0=0$ convention. Exponential weights give the
following identity for every comparator $q_i$:
\begin{align*}
 \eta\ip{q_i-x_i^t}{v_i^t}
 ={}&D(q_i\Vert y_i^t)-D(q_i\Vert y_i^{t+1})
 -D(y_i^{t+1}\Vert x_i^t)-D(x_i^t\Vert y_i^t)\\
 &+\eta\ip{y_i^{t+1}-x_i^t}{v_i^t-v_i^{t-1}}.
\end{align*}
To verify it, apply the log-ratio identities for
$y_i^{t+1}\propto y_i^t e^{\eta v_i^t}$ and
$x_i^t\propto y_i^t e^{\eta v_i^{t-1}}$, and add them.
Pinsker's inequality $D(z\Vert w)\geq\norm{z-w}_1^2/2$ and
$\eta a d\leq a^2/4+\eta^2d^2$ imply, after summing in time and maximizing
over $q_i$,
\begin{equation}\label{eq:rvu-l1}
 R_i(T)\leq\frac{\log m_i}{\eta}
 +\eta\sum_{t\leq T}d_{it}^2
 -\frac1{4\eta}\sum_{t\leq T}(a_{it}^2+b_{it}^2).
\end{equation}
This is the standard optimistic mirror-descent calculation, with constants
given here explicitly; predictable-sequence methods are developed in
\cite{rs}.

Put $E_T=\sum_{i,t\leq T}(a_{it}^2+b_{it}^2)$ and
$\delta_{jt}=\norm{x_j^t-x_j^{t-1}}_1$ for $t\geq2$.
Bounded payoffs and total variation of product distributions give
\[
 d_{it}\leq\frac12\sum_{j\ne i}\delta_{jt}\quad(t\geq2),
 \qquad d_{i1}\leq1.
\]
Indeed, changing one opponent's distribution changes the expectation of a
$[0,1]$-valued payoff by at most half its $\ell_1$ change, and replacing
opponents one at a time proves the displayed estimate.
Also $\delta_{jt}\leq b_{jt}+a_{j,t-1}$, so
$\sum_{j,t=2}^T\delta_{jt}^2\leq2E_T$. Consequently,
\begin{equation}\label{eq:all-feedback}
 \sum_{i,t\leq T}d_{it}^2
 \leq N+\frac{(N-1)^2}{2}E_T,
 \qquad
 \sum_{t\leq T}d_{it}^2
 \leq1+\frac{N-1}{2}E_T.
\end{equation}

The external-regret maximum is at least the regret against each $p_i$.
Summing \eqref{eq:aggregate-condition} over rounds therefore proves
$\sum_iR_i(T)\geq0$. Summing \eqref{eq:rvu-l1} and applying
\eqref{eq:all-feedback},
\[
 0\leq\sum_iR_i(T)
 \leq\frac L\eta+\eta N+
 \left(\frac{\eta(N-1)^2}{2}-\frac1{4\eta}\right)E_T.
\]
The step-size restriction makes the coefficient at most $-1/(8\eta)$.
Hence $E_T\leq8L+8\eta^2N$. Substitute this into the individual bound in
\eqref{eq:all-feedback}, then drop the negative term in
\eqref{eq:rvu-l1}, to obtain \eqref{eq:aggregate-bound}.
At the stated explicit step size, its right side is at most
\[
 (4N-2)\log K+\frac1{2(N-1)}+\frac{N}{2(N-1)^2}
 \leq(4N-2)\log K+\frac32.
\]
\end{proof}

\begin{corollary}\label{cor:zero-sum}
In every two-player constant-sum game with each payoff in $[0,1]$,
plain Optimistic Hedge with $\eta=1/2$ satisfies
$R_i(T)\leq6\log K+3/2$ for all $i,T$.
\end{corollary}

\begin{proof}
Let $u_1+u_2=c$, and let $p_1,p_2$ be minimax strategies for $u_1$,
with value $V$. For every $x_1,x_2$,
$u_1(p_1,x_2)\geq V\geq u_1(x_1,p_2)$. Therefore
\[
 \sum_{i=1}^2[u_i(p_i,x_{-i})-u_i(x)]
 =u_1(p_1,x_2)-u_1(x_1,p_2)\geq0.
\]
Apply Theorem~\ref{thm:aggregate} with $N=2$.
\end{proof}
The qualitative constant-regret phenomenon in this class was already known.
Tsuchiya studies substantially sharper action-count dependence and constants
\cite{tsuchiya}; the corollary is a directly checked benchmark for the exact
initialization in this problem.

\subsection*{4. A lower bound that every proposed step-size rule must respect}

\begin{proposition}\label{prop:lower}
For every $N,K\geq2$ and every fixed $\eta>0$, there is an allowed game in
which
\[
 \sup_{T\geq1}R_i(T)\geq\frac{\log K}{\eta}-\frac32.
\]
Thus any uniform bound for a proposed rule $\eta(N,K)$ must be at least
this quantity (as well as nonnegative).
\end{proposition}

\begin{proof}
Give player $i$ exactly $K$ actions, with payoff one for one action and zero
for all the others, independently of the opponents. Give all other players
arbitrary allowed action sets and constant payoffs. Put $q=K-1$.
At $t=1$ the regret increment is $q/K$. At $t\geq2$,
\eqref{eq:telescoped} gives the good action probability
$e^{\eta t}/(e^{\eta t}+q)$, so
\[
 \sup_T R_i(T)=\frac qK+\sum_{t=2}^\infty\frac q{e^{\eta t}+q}.
\]
The summand is decreasing, and comparison with its integral gives
\[
 \sum_{t=2}^\infty\frac q{e^{\eta t}+q}
 \geq\int_2^\infty\frac q{e^{\eta s}+q}\,ds
 =\frac1\eta\log(1+q e^{-2\eta})
 \geq\frac{\log K}{\eta}-2.
\]
Since $q/K\geq1/2$, the claimed bound follows.
\end{proof}
This example rules out making an extremely small step size costless. It
does not disprove a polynomial bound: that would require controlling all
possible dimension-dependent step-size rules.

\subsection*{5. Making the finite-recurrence argument quantitative}
The next theorem replaces all constants after the recurrence step by
explicit formulas. It identifies a finite set of quantitative data that
would suffice to settle the general problem.

Use losses $\ell_i^t=1-v_i^t$ for $t\geq1$, with $\ell_i^0=0$.
This produces the same normalized update as \eqref{eq:update}; the
action-independent first-step factor cancels. For $a<b$ define
\[
 h_{i,ab}^t=\sqrt{x_{ia}^tx_{ib}^t}
 (\ell_{ia}^{t-1}-\ell_{ib}^{t-1}),\qquad
 \Delta z^t=z^{t+1}-z^t.
\]
In particular $h_i^1=0$, and the pairwise variance identity gives
$\norm{h_i^t}_2^2=\Var_{x_i^t}(\ell_i^{t-1})\leq1/4$ for $t\geq2$.

\begin{theorem}[Quantitative recurrence criterion]\label{thm:certificate}
Suppose numerical values $m\geq1$, $A,M\geq0$, and $\eta_0>0$ have been
proved to work for every player and every game in a specified class, in the
following sense: for all $0<\eta\leq\eta_0$ and all $t\geq1$,
\begin{align}
 \Delta^m h_i^{t+m}
 &=\sum_{k=0}^{m-1}c_{ikt}\eta^{m-k}\Delta^k h_i^{t+k},
 & |c_{ikt}|&\leq A,\label{eq:certificate-recurrence}\\
 \norm{\Delta^k h_i^{t+k}}_2&\leq M\eta^k,
 &0\leq k&\leq m.\label{eq:certificate-smooth}
\end{align}
The coefficients may depend on the player, game, time, and step size.
Define explicit quantities
\[
 P=2m\max\{A,1\},\quad
 E=3^mM\sqrt{8m},\quad
 D=16mM^2,\quad
 Q=2(2E)^{2/m},\quad C=2P^2+Q,
\]
and
\[
 \bar\eta=\min\!\left\{\eta_0,\frac1{12},
          \frac1{\sqrt{6(C+32)}}\right\}.
\]
For every $0<\eta\leq\bar\eta$, every player, and every horizon,
\begin{equation}\label{eq:certificate-bound}
 R_i(T)\leq\frac{\log m_i}{\eta}
             +\eta\left(D+Q+\frac1{24}\right).
\end{equation}
Thus supplied uniform recurrence data produce explicit functions
$\eta=\bar\eta$ and
$B=(\log K)/\bar\eta+\bar\eta(D+Q+1/24)$ for that class.
\end{theorem}

\subsubsection*{5.1. An explicit cutoff lemma}

\begin{lemma}\label{lem:cutoff}
Let $0<\eta\leq1$ and let a Euclidean-valued sequence $(u^t)_{t\geq1}$
satisfy \eqref{eq:certificate-recurrence}--\eqref{eq:certificate-smooth}
with $h_i$ replaced by $u$. For every $T\geq1$, there is a finitely
supported sequence $y$ on $\mathbb Z$ such that, writing
$U_j=\norm{\Delta^j y}_{\ell_2}$,
\begin{align}
 U_0^2&\leq\sum_{t=1}^{T+1}\norm{u^t}_2^2,\label{eq:cutoff-zero}\\
 \sum_{t=1}^T\norm{\Delta u^t}_2^2&\leq U_1^2+D,\label{eq:cutoff-one}\\
 U_m&\leq A\sum_{k=0}^{m-1}\eta^{m-k}U_k
              +E\eta^{m-1/2}.\label{eq:cutoff-high}
\end{align}
Here $D=16mM^2$ and $E=3^mM\sqrt{8m}$.
\end{lemma}

\begin{proof}
Set $J=\lceil1/\eta\rceil$ and $w=m(J-1)$. If
$T\leq2w+6m+1$, choose $y=0$. The order-zero bound gives
$\norm{\Delta u^1}\leq2M$, and the order-one bound gives
$\norm{\Delta u^t}\leq M\eta$ for $t\geq2$. Since $w\leq m/\eta$,
\[
 \sum_{t=1}^T\norm{\Delta u^t}^2
 \leq4M^2+TM^2\eta^2\leq(8m+5)M^2\leq D.
\]
The other conclusions are immediate.

For a longer horizon let $a=3m+1$, $b=T+1-3m$. Let $b_J$ be the
probability mass $1/J$ on $\{0,\ldots,J-1\}$ and zero elsewhere, and let
$q=b_J^{*m}$ be its $m$-fold discrete convolution. It is nonnegative,
has sum one, and is supported on $\{0,\ldots,w\}$. Define
\[
 \chi=1_{\{a,\ldots,b-w\}}*q,\qquad
 y^t=\chi^t u^t\ (t\geq1),\qquad y^t=0\ (t\leq0).
\]
Then $0\leq\chi\leq1$, its support is contained in $[a,b]$, and
$\chi=1$ on $[a+w,b-w]$. This proves \eqref{eq:cutoff-zero}.
Because $\norm{\Delta b_J}_{\ell_1}=2/J$, convolution gives
\[
 \sup_t|\Delta^r\chi^t|\leq(2/J)^r\leq(2\eta)^r,
 \qquad 0\leq r\leq m.
\]
Let $\mathcal B$ be the integer starting indices whose stencil
$\{t,\ldots,t+2m\}$ meets $[a,a+w]$ or $[b-w,b]$.
Then
\[
 |\mathcal B|\leq2w+4m+2\leq8m/\eta,
 \qquad \min\mathcal B\geq a-2m=m+1.
\]
For $t\in\mathcal B$, all needed differences of $u$ therefore begin in
the range where \eqref{eq:certificate-smooth} applies. The product rule
\[
 \Delta^j(\chi u)^s
 =\sum_{r=0}^j\binom jr(\Delta^r\chi)^s
                         (\Delta^{j-r}u)^{s+r}
\]
implies $\norm{\Delta^j y^s}\leq3^jM\eta^j$ on the stencils in
question, for $0\leq j\leq m$.

Outside $\mathcal B$, a stencil has all cutoff values one or all zero.
Use the original recurrence coefficients on the former stencils and zero
coefficients everywhere else. The residual in the recurrence for $y$ is
then supported in $\mathcal B$ and bounded pointwise by $3^mM\eta^m$.
Its $\ell_2$ norm is at most $E\eta^{m-1/2}$. The triangle inequality and
invariance of $\ell_2$ norms under integer shifts prove
\eqref{eq:cutoff-high}.

Finally, $\Delta y^t=\Delta u^t$ whenever
$a+w\leq t\leq b-w-1$. At most $2w+6m+1\leq(8m+1)/\eta$ other
indices $t\in\{2,\ldots,T\}$ remain. Their contribution is at most
$(8m+1)M^2\eta$, and the initial difference contributes at most $4M^2$.
This gives \eqref{eq:cutoff-one} with the stated $D$.
\end{proof}

\subsubsection*{5.2. Interpolation and numerical constants}
For finite-support Hilbert-valued sequences, Parseval's identity and
H\"older's inequality give
\begin{equation}\label{eq:interpolation}
 U_j\leq U_0^{1-j/m}U_m^{j/m}\qquad(0\leq j\leq m).
\end{equation}
Indeed the Fourier multiplier for $\Delta$ has squared modulus
$|e^{i\theta}-1|^2$; applying H\"older to its $j$th power proves the
claim after integration. This proof also covers vector-valued sequences.

Combine \eqref{eq:cutoff-high} with \eqref{eq:interpolation}. We claim
\begin{equation}\label{eq:highest-explicit}
 U_m\leq P^m\eta^m U_0+2E\eta^{m-1/2}.
\end{equation}
If $U_0=0$ all terms vanish. Otherwise, either
$U_m\leq P^m\eta^m U_0$, or put
$r=(U_m/(\eta^mU_0))^{1/m}>P\geq1$. In the latter case the sum on the
right side of \eqref{eq:cutoff-high} is at most
$AU_m\sum_{k<m}r^{k-m}\leq(Am/r)U_m\leq U_m/2$.
Absorption proves \eqref{eq:highest-explicit} in either case.

For $m\geq1$ and nonnegative $z_1,z_2$,
$(z_1+z_2)^{2/m}\leq2(z_1^{2/m}+z_2^{2/m})$.
Using \eqref{eq:interpolation}, \eqref{eq:highest-explicit}, and weighted
Young's inequality gives
\begin{align*}
 U_1^2
 &\leq2P^2\eta^2U_0^2
       +2(2E)^{2/m}(\eta^2U_0^2)^{1-1/m}\eta^{1/m}\\
 &\leq C\eta^2U_0^2+Q\eta,
\end{align*}
with the stated $Q$ and $C$. At $m=1$ this is read directly as
$U_1^2\leq2P^2\eta^2U_0^2+Q\eta$; thus no ambiguous $0^0$
convention is needed. Lemma~\ref{lem:cutoff} proves
\begin{equation}\label{eq:temporal-explicit}
 \sum_{t=1}^T\norm{\Delta h_i^t}_2^2
 \leq C\eta^2\sum_{t=1}^{T+1}\norm{h_i^t}_2^2+D+Q\eta.
\end{equation}

\subsubsection*{5.3. Completing the regret calculation}
For completeness, the one-player variance inequalities used next are
proved here with numerical constants. They are the variance form of the
optimistic potential calculation used in Ha \cite{ha}.
Define $\Phi(z)=-(1/\eta)\log\sum_a e^{-\eta z_a}$ and let
$D_{-\Phi}$ denote the Bregman divergence of the convex function
$-\Phi$. If $p=\nabla\Phi(z)$ and the coordinate span of $w$ is at most
$r$, its integral Hessian formula gives
\[
 \frac\eta2 e^{-\eta r}\Var_p(w)
 \leq D_{-\Phi}(z+w,z)
 \leq\frac\eta2 e^{\eta r}\Var_p(w).
\]
To check this, along the segment $z+sw$ each probability is between
$e^{-\eta r}$ and $e^{\eta r}$ times its initial value. Apply these bounds
to $\Var_p(w)=\min_c\sum_a p_a(w_a-c)^2$ and integrate $\eta(1-s)$
over $0\leq s\leq1$.

Let $L^t=\sum_{s\leq t}\ell^s$ and $q^t=L^{t-1}+\ell^{t-1}$.
Then $x^t=\nabla\Phi(q^t)$, and direct expansion gives
\[
 D_{-\Phi}(L^t,q^t)-D_{-\Phi}(L^{t-1},q^t)
 =\ip{x^t}{\ell^t}-\Phi(L^t)+\Phi(L^{t-1}).
\]
Sum this identity and compare $\Phi(L^T)$ with $\min_aL_a^T$.
The coordinate spans of $\ell^t-\ell^{t-1}$ and $-\ell^{t-1}$ are at
most two and one respectively. Thus, for $\eta\leq1/12$,
\begin{equation}\label{eq:variance-rvu}
 R_i(T)\leq\frac{\log m_i}{\eta}
 +\frac{2\eta}3\sum_{t=1}^T\Var_{x_i^t}(\ell_i^t-\ell_i^{t-1})
 -\frac\eta3\sum_{t=1}^T\norm{h_i^t}_2^2.
\end{equation}
Here $e^{2\eta}/2\leq2/3$ and $e^{-\eta}/2\geq1/3$ on the stated
interval, which also follows from $e^z\leq1/(1-z)$ for $0\leq z<1$.

Write $w_{ab}^t=\sqrt{x_{ia}^tx_{ib}^t}$. The coordinate span of
$2\ell_i^t-\ell_i^{t-1}$ is at most three, so
$e^{-3\eta}\leq w_{ab}^{t+1}/w_{ab}^t\leq e^{3\eta}$.
For $\eta\leq1/12$, $e^{3\eta}-1\leq4\eta$. Hence
\begin{align*}
 &w_{ab}^t\bigl[(\ell_{ia}^t-\ell_{ib}^t)
                   -(\ell_{ia}^{t-1}-\ell_{ib}^{t-1})\bigr]\\
 &\qquad=(h_{i,ab}^{t+1}-h_{i,ab}^t)
       +\left(\frac{w_{ab}^t}{w_{ab}^{t+1}}-1\right)h_{i,ab}^{t+1}.
\end{align*}
Square with $(z+w)^2\leq(3/2)z^2+3w^2$, then sum over pairs and time.
The pairwise variance identity and the terminal bound
$\norm{h_i^{T+1}}^2\leq1/4$ give
\begin{equation}\label{eq:weighted-gap-explicit}
 \sum_{t=1}^T\Var_{x_i^t}(\ell_i^t-\ell_i^{t-1})
 \leq\frac32\sum_{t=1}^T\norm{\Delta h_i^t}_2^2
 +48\eta^2\left(\sum_{t=1}^T\norm{h_i^t}_2^2+\frac14\right).
\end{equation}
Insert \eqref{eq:temporal-explicit} and
\eqref{eq:weighted-gap-explicit} into \eqref{eq:variance-rvu}.
Writing $H=\sum_{t=1}^T\norm{h_i^t}^2$, we obtain
\[
 R_i(T)\leq\frac{\log m_i}{\eta}+\eta D+Q\eta^2
 +\left((C+32)\eta^3-\frac\eta3\right)H
 +\frac{C+32}{4}\eta^3.
\]
The definition of $\bar\eta$ makes the coefficient of $H$ at most
$-\eta/6$ and the last term at most $\eta/24$.
Drop the nonpositive term and use $\eta\leq1$ on $Q\eta^2$.
This proves Theorem~\ref{thm:certificate}.

\subsection*{6. Conditional computability for all finite games}
Here is a precise conditional way to make the outstanding recurrence search
effective. Its hypothesis is a decision procedure for the full first-order
theory of
\[
 \mathbb R_{\exp}=(\mathbb R;0,1,+,\cdot,<,\exp).
\]
Only sentences with rational coefficients are needed. This hypothesis is
not supplied by real analyticity or by Noetherianity alone. The real form
of Schanuel's conjecture implies such decidability; Berarducci and
Gallinaro state the precise implication and give a new conditional
axiomatization in \cite{bg}. The theorem below needs only the decision
procedure, not the stronger number-theoretic conjecture as a separate
premise.

\begin{theorem}[Conditional effective search]\label{thm:conditional-search}
Assume that the first-order theory of $\mathbb R_{\exp}$ is decidable.
Then, using the qualitative finite-generation theorem and analytic state
construction in Ha \cite[Theorem 5.5 and Lemmas 5.6--5.8]{ha}, there is a
terminating algorithm which, on input $(N,K)$, computes a positive common
step size and a finite regret bound for the exact update
\eqref{eq:update}, uniformly over all allowed games and horizons.
This is a conditional answer to the computability part of prb_bc572aa70a8748a3b489fa9941f9307e. It
does not imply the polynomial subquestion.
\end{theorem}

\begin{proof}
Fix an action-count vector $\mathbf d=(d_1,\ldots,d_N)$ with
$1\leq d_i\leq K$. The parameter space of loss tables is
$\mathcal G=[0,1]^{N\prod_j d_j}$. The real state space is the compact
semialgebraic set
\[
 \mathcal S=\{(r,z): r_{ia}\geq0,\ \sum_a r_{ia}^2=1,
                          \ 0\leq z_{ia}\leq1\text{ for all }i,a\}.
\]
Here $r_{ia}$ is a square-root probability and $z_i$ stores the preceding
loss vector. For real $r$ near the product of unit spheres, set
\[
 p_{ia}(r)=\frac{r_{ia}^2}{\sum_b r_{ib}^2},\qquad
 L_{g,ia}(r)=\sum_{a_{-i}}g_i(a,a_{-i})\prod_{j\ne i}p_{j,a_j}(r).
\]
With $s=(r,z)$, define $G(s;g,\eta)=(F(s;g,\eta),L_g(r))$, where
\[
 F_{ia}(s;g,\eta)=
 \frac{\left(\sum_b r_{ib}^2\right)^{1/2}r_{ia}
       \exp\{-\eta(2L_{g,ia}(r)-z_{ia})/2\}}
 {\left(\sum_b r_{ib}^2
       \exp\{-\eta(2L_{g,ib}(r)-z_{ib})\}\right)^{1/2}}.
\]
The table $g$ and step size $\eta$ remain fixed under composition of $G$.
For real inputs in $\mathcal S$ every denominator is positive, and $G$
maps $\mathcal S$ into itself. The extra numerator norm is one on the
state space; it ensures the ambient identity
$G(r,z;g,0)=(r,L_g(r))$ and hence $G^2(s;g,0)=G(s;g,0)$.

Let $H_0(s)$ be the vector formed by stacking, for all players, the
coordinates $r_{ia}r_{ib}(z_{ia}-z_{ib})$, $a<b$. For $\eta>0$, define
the finite expression
\begin{equation}\label{eq:explicit-H}
 H_k(s;g,\eta)=\eta^{-k}
 \sum_{j=0}^k(-1)^{k-j}\binom kj
          H_0\bigl(G^{k+j}(s;g,\eta)\bigr).
\end{equation}
Each $H_k$ uses at most $2k$ iterates and finitely many field operations,
positive square roots, and exponentials. Along the actual trajectory
$s^t=(\sqrt{x^t},\ell^{t-1})$,
\[
 H_k(s^t;g,\eta)=\eta^{-k}\Delta^k h^{t+k},
\]
where $h$ stacks the players' weighted gaps. In particular the order-one
expression controls $\Delta h^{t+1}$, not $\Delta h^t$ at $t=1$.
The exceptional initial difference remains covered by the cutoff proof
above.

For positive integers $(m,A,M,q)$ form the following first-order sentence:
for every $s\in\mathcal S$, every $g\in\mathcal G$, and every
$0<\eta\leq1/q$,
\begin{equation}\label{eq:fo-certificate}
 \begin{gathered}
 \exists c_0,\ldots,c_{m-1}\in\mathbb R:
       \quad |c_k|\leq A\ (k<m),\qquad
       H_m(s;g,\eta)=\sum_{k<m}c_k H_k(s;g,\eta),\\
 \norm{H_k(s;g,\eta)}_2^2\leq M^2
                   \qquad(0\leq k\leq m).
 \end{gathered}
\end{equation}
The universal quantifiers have the indicated antecedents; coefficients
$c_k$ may depend on $s,g,\eta$. The norm squares and vector identities
are finite sums and coordinate conjunctions. Replace every positive
square root by an auxiliary variable $b>0$ with its defining equation
$b^2=a$; replace divisions by their defining equations with nonzero
denominators. This effectively converts \eqref{eq:fo-certificate} to a
sentence in the language of $\mathbb R_{\exp}$ with rational
coefficients. No quantification over functions or over infinitely many
iterates occurs in a candidate sentence.

We next verify that at least one candidate is true. The compact set for
the analytic argument is
$\mathcal S\times\mathcal G\times\{0\}$, including the stored-loss
coordinate independently of the current expected loss. Although
\eqref{eq:explicit-H} displays division by $\eta$, each fixed $H_k$
extends analytically through zero. Indeed, equivalently,
\[
 H_{k+1}=\bigl(H_k\circ G^2-H_k\circ G\bigr)/\eta.
\]
The numerator vanishes at $\eta=0$ by the idempotence of $G(\cdot;g,0)$.
Induction proves analytic divisibility, with a neighborhood allowed to
shrink at each finite order. Ha's finite-generation theorem applies to the
stacked vectors $H_0,H_1,\ldots$ just as it applies to an individual
player's vector. It yields a finite $m$ and analytic scalar coefficients
expressing $H_m$ in terms of $H_0,\ldots,H_{m-1}$ on a neighborhood of
that compact set. Compactness then supplies a positive common step-size
interval and finite upper bounds for the finitely many coefficients and
$H_k$ involved. Choosing sufficiently large integers $A,M,q$ proves
that \eqref{eq:fo-certificate} is true for at least one tuple. This
existence argument does not itself compute the tuple.

Now enumerate all positive-integer tuples $(m,A,M,q)$ in increasing sum,
breaking ties lexicographically. Apply the assumed decision procedure to
each associated sentence. Each test terminates. Since a true candidate
exists, the search eventually returns one. Its guarantees hold on the
actual trajectory at every round by invariance of $\mathcal S$.
The stacked norm bounds imply the corresponding bounds for each player's
subvector, so Theorem~\ref{thm:certificate} applies with
$\eta_0=1/q$.

There are only finitely many action-count vectors
$\mathbf d\in\{1,\ldots,K\}^N$. Run this terminating search for each.
For each vector compute $\bar\eta_{\mathbf d}$ and its constants
$D_{\mathbf d},Q_{\mathbf d}$ from Theorem~\ref{thm:certificate}.
Choose
\[
 \eta(N,K)=\min_{\mathbf d}\bar\eta_{\mathbf d},\qquad
 B(N,K)=\frac{\log K}{\eta(N,K)}
     +\eta(N,K)\max_{\mathbf d}
                 \left(D_{\mathbf d}+Q_{\mathbf d}+\frac1{24}\right).
\]
These are positive computable real quantities, and all their nonlogarithmic
operations are algebraic operations on the returned integers. If a
rational step size and integer bound are preferred, choose a smaller
positive rational step size by algebraic comparison and an integer upper
bound using certified rational approximations to $\log K$.
Theorem~\ref{thm:certificate} proves the required horizon-uniform bound.
\end{proof}

\subsection*{7. What is proved and what is still missing}
Theorems~\ref{thm:dominance} and \ref{thm:aggregate} are unconditional
on their stated classes. Theorem~\ref{thm:certificate} is a proved
implication with numerical constants, not an assumption-free solution of
prb_bc572aa70a8748a3b489fa9941f9307e. It reduces the outstanding quantitative task to producing,
from $(N,K)$ alone, valid numbers $(m,A,M,\eta_0)$ in
\eqref{eq:certificate-recurrence}--\eqref{eq:certificate-smooth}, or
an alternative argument of comparable force.

Ha supplies existence through an analytic module near the compact set of
states and game tables at zero step size. In his notation the module
generators are the analytic scaled differences $H_{i,k}$, and finite
generation produces a relation $H_{i,m}=\sum_{k<m}c_kH_{i,k}$.
Merely saying ``enumerate possible parameters'' does not provide a
terminating verification procedure for that analytic identity. Section 6
supplies the missing verifier under the explicitly additional decidability
hypothesis. Compactness alone does not compute a neighborhood radius, and
no unconditional certificate-search algorithm is established here.

There is a second unresolved question: even an extremely large computable
bound would not establish polynomial dependence on $N$ and $\log K$.
Conversely, the static-payoff lower bound in
Proposition~\ref{prop:lower} constrains a proposed step size but does not
exclude all polynomially suitable rules. The full answer and the polynomial
subquestion therefore both remain \textbf{UNRESOLVED} unconditionally.

\subsection*{8. Checks performed}
The proofs above retain the original initialization and common fixed step
size; handle $m_i=1$; distinguish cumulative regret from its absolute
value; and control the exceptional first difference separately in the
cutoff construction. The bounds are independent of the horizon and, on
their stated classes, of the payoff table. No numerical trajectory is
used as proof. The primary Ha manuscript was inspected in its exact
version \texttt{2609.22839v1}; the source theorem and the explicitly
stated quantitative gap match the catalogue formulation. The displayed
special cases and reductions have been checked algebraically within this
contribution, without a claim of formal verification or novelty.

\section*{Models used}
OpenAI GPT-6 Astra Pro: mathematical derivation, adversarial checking,
primary-source inspection, and drafting. Exact internal build identifier
and formal reasoning configuration: unknown. The work was prepared by a
parallel research agent within the same assistant session, followed by
separate-agent review of the quantitative recurrence proof and conditional
certificate search, followed by root review.
No external referee approval is claimed.

\section*{Human contributors and roles}
The user requested the problem-by-problem solution attempts and supplied
the problem package. No independent human proof audit is claimed.
The original problem suggestions are credited in the supplied package to
Rida Laraki. The results cited below retain their original authors' credit.

\begin{thebibliography}{9}
\raggedright
\bibitem{ha}
Junsoo Ha. \emph{A Horizon-Independent Regret Bound for Optimistic Hedge in
General-Sum Games}. arXiv:2609.22839v1, 19 September 2026.
Equation (5), Theorem 4.1, text following Corollary 4.2, Sections 5.3--5.5,
and Appendices B--D inspected.
\url{https://arxiv.org/html/2609.22839v1}.

\bibitem{rs}
Alexander Rakhlin and Karthik Sridharan.
\emph{Online Learning with Predictable Sequences}.
Proceedings of the 26th Annual Conference on Learning Theory,
PMLR 30:993--1019, 2013. Official publication metadata and abstract checked;
the identities needed here are proved in the text.
\url{https://proceedings.mlr.press/v30/Rakhlin13.html}.

\bibitem{tsuchiya}
Taira Tsuchiya. \emph{Tight Regret Upper and Lower Bounds for Optimistic
Hedge in Two-Player Zero-Sum Games}. arXiv:2510.11691v1,
13 October 2025. Primary arXiv metadata and abstract checked; cited for
the existing zero-sum literature, not as a proof dependency.
\url{https://arxiv.org/abs/2510.11691}.

\bibitem{bg}
Alessandro Berarducci and Francesco Gallinaro.
\emph{On the elementary theory of the real exponential field}.
arXiv:2603.08365v1, 9 March 2026. Abstract and Introduction inspected for
the implication from the real form of Schanuel's conjecture to
decidability, and for the distinction between model completeness and
completeness of the restricted-exponential axioms.
\url{https://arxiv.org/html/2603.08365v1}.
\end{thebibliography}
\end{document}
